\[ %%%%%%%%%%%%%%%%%%%%%%%%%%%%%%%%%%%%%%%%%%%%%%%%%%%%%%%%%%%%%%%%%%%%%%%%%%%%%%%% \subsection{Digital Twins and Scientific Foundation Models} \label{sec:digital_twins} %%%%%%%%%%%%%%%%%%%%%%%%%%%%%%%%%%%%%%%%%%%%%%%%%%%%%%%%%%%%%%%%%%%%%%%%%%%%%%%% Digital twins have emerged as a transformative paradigm for representing, monitoring, and predicting the behavior of complex physical systems through continuously updated computational models. Unlike traditional numerical simulations that operate as isolated computational experiments, digital twins aim to establish persistent connections between physical assets and their virtual representations by integrating observations, simulation models, machine learning algorithms, and real-time data streams. A general digital twin architecture consists of several interacting components: a physical system, a sensing and data acquisition layer, a computational model, a prediction and optimization engine, and a feedback mechanism that enables continuous synchronization between the physical and virtual domains. Such architectures have been increasingly applied in aerospace engineering, manufacturing, energy systems, healthcare, environmental monitoring, and infrastructure management. However, the effectiveness of digital twins strongly depends on the ability of their computational models to accurately represent complex system dynamics over extended time horizons. Conventional physics-based simulations provide high interpretability and physical consistency but are often computationally expensive, limiting their application in real-time environments. Conversely, purely data-driven machine learning models provide rapid prediction capabilities but may suffer from limited extrapolation ability and violation of physical constraints. This challenge has motivated the integration of scientific machine learning methods into digital twin architectures. Physics-informed neural networks, neural operators, graph neural simulators, and differentiable physics models have been explored as potential computational engines for next-generation digital twins. These approaches aim to combine the predictive efficiency of machine learning with the reliability of physics-based modeling. Graph-based approaches are particularly attractive for digital twins because many physical systems naturally consist of interacting components. Infrastructure networks, biological systems, mechanical assemblies, and environmental processes can be represented as dynamic graphs where nodes correspond to physical entities and edges represent interactions or information exchange. Such representations provide a natural foundation for adaptive digital twins capable of modeling complex relational structures. Despite significant progress, current digital twin implementations still face several fundamental limitations. Many existing systems rely on collections of specialized models designed for individual applications rather than general computational architectures capable of transferring knowledge across different physical domains. Furthermore, maintaining long-term predictive accuracy while continuously integrating new observations remains a major challenge. Recent developments in artificial intelligence have introduced the concept of scientific foundation models: large-scale computational models designed to learn general representations of scientific data, physical processes, and mathematical structures. Inspired by the success of foundation models in natural language and computer vision, scientific foundation models aim to provide reusable representations that can accelerate discovery and simulation across multiple scientific domains. Examples of this emerging direction include large-scale models for weather prediction, molecular representation learning, climate simulation, and universal neural operators. These systems attempt to move beyond task-specific prediction toward general-purpose scientific reasoning and simulation capabilities. However, scientific foundation models introduce additional requirements that distinguish them from conventional AI systems. A useful scientific foundation model must preserve physical consistency, support extrapolation beyond training distributions, represent multi-scale interactions, quantify uncertainty, and remain interpretable with respect to underlying scientific principles. From this perspective, physics-guided architectures provide an important foundation for future scientific foundation models. Rather than relying exclusively on statistical correlations extracted from large datasets, physics-guided models incorporate structural knowledge about the systems they represent. This enables more reliable generalization across physical regimes and improves robustness when data are limited. The PHYSGEN framework contributes to this emerging direction by introducing a physics-guided graph architecture designed for learning and predicting the evolution of complex dynamical systems. By combining graph-based representations, differentiable physical constraints, and adaptive stability mechanisms, PHYSGEN provides a computational framework that can serve as a building block for future physics-guided digital twins and scientific foundation models. In this context, PHYSGEN represents a step toward a new generation of scientific artificial intelligence systems in which learning, simulation, and physical reasoning are integrated into a unified computational paradigm. \] 
